\begin{center}
\refstepcounter{table}\label{tab:rectilinear_baseline}
\begin{minipage}{\columnwidth}
\small\textbf{Table \thetable.} Rectilinear no-force baseline for present-day Galactocentric phase space.
\end{minipage}
\vspace{0.4ex}

\scriptsize
\setlength{\tabcolsep}{2.5pt}
\begin{tabular}{lrrrrr}
\hline\hline
Sample & $N$ & $\tilde d_{\rm lin}$ & $\widetilde{|t_{\rm lin}|}$ & MC $N(d<0.5{\rm\,kpc})$ & MC $N(d<1{\rm\,kpc})$ \\
\hline
Tier~A+B & 276 & 6.91 & 260 & 0 [0, 1] & 2 [1, 4] \\
Gold A+B & 226 & 6.91 & 258 & 0 [0, 1] & 2 [1, 3] \\
Tier~A+B+C & 621 & 6.81 & 240 & 1 [0, 2] & 3 [2, 5] \\
\hline
\end{tabular}
\vspace{0.5ex}

\begin{minipage}{0.95\columnwidth}
\footnotesize
\textit{Note.} The baseline uses $\mathbf{r}(t)=\mathbf{r}_0+\mathbf{v}_0t$ in the adopted Galactocentric frame, with no Galactic potential.  Positive $t_{\rm lin}$ is future and negative $t_{\rm lin}$ is past; $d_{\rm lin}$ is the minimum straight-line Galactocentric distance.  Point-estimate medians of $d_{\rm lin}$ are in kpc and $|t_{\rm lin}|$ in Myr.  MC count intervals are 16th--84th percentiles from 5{,}000 present-day Gaia phase-space realisations per star.
\end{minipage}
\end{center}
